%!LW recipe=pdfLaTeXmk

\documentclass[scheme=plain]{beamer}
\usepackage[T1]{fontenc}
\usepackage{fontawesome5}
\usepackage{mathtools}
\usepackage{tikz}
\usepackage{booktabs}
\usepackage{caption}
\usepackage{outlines}
\usepackage{graphicx}
\usepackage{float}
\usepackage{amsthm}
\usepackage{tabularray}
\usepackage{minted}
\usepackage{hyperref}
\usepackage{cleveref}
\usepackage{url}
\usepackage{xspace}
\usepackage{academicons}
\usepackage{pgfplots}
\usepackage{pgfgantt}
\usepackage{qrcode}
\usepackage{ninecolors}
\usepackage{soul}
\usepackage{tgheros}
\usepackage[style=numeric, backend=biber, url=false, doi=false, isbn=false]{biblatex}
\newmintinline[code]{html}{fontsize=\small}
\addbibresource{custom.bib}
\usepackage{graphicx}
\usetheme[
    progressbar=frametitle,
    numbering=fraction,
    subsectionpage=progressbar,
    titleformat title=smallcaps,
    titleformat subtitle=smallcaps,
    titleformat section=smallcaps,
    titleformat frame=smallcaps]{moloch}

    
\setbeamercolor{palette primary}{bg=gray1}
\setbeamercolor{frametitle}{bg=black,fg=white}
\setbeamercolor{progress bar}{fg=black}
\setbeamerfont{footnote}{size=\tiny}
\renewcommand{\cite}{\footfullcite}


\definecolor{ForestGreen}{RGB}{34,139,34}
\newcommand{\highlightRelevantValue}[1]{%
  \ifdim #1pt > 0pt
    {\color{ForestGreen}{#1\%}}%
  \else
    {\color{red}{#1\%}}%
  \fi
}

\usetikzlibrary{calc}
\setcounter{tocdepth}{1}

\makeatletter
\newlength{\frametitle@padding}
\setlength{\frametitle@padding}{2.2ex}
\newcommand{\frametitlestrut@start}{
  \rule{0pt}{\frametitle@padding +%
    \totalheightof{%
      \ifcsdef{frametitleformat}{\frametitleformat X}{X}%
    }%
  }%
}

\newcommand{\frametitlestrut@end}{
  \rule[-\frametitle@padding]{0pt}{\frametitle@padding}
}

\setbeamertemplate{frametitle}{
    \nointerlineskip%
    \begin{beamercolorbox}[%
        wd=\paperwidth,%
        sep=0pt,%
        leftskip=\frametitle@padding,%
        rightskip=\frametitle@padding,%
      ]{frametitle}%
    \frametitlestrut@start%
    \insertframetitle%
    \hfill%
    \nolinebreak%
    \frametitlestrut@end%
    \end{beamercolorbox}%
    \begin{tikzpicture}[remember picture,overlay]
        \coordinate (logo) at ([xshift=-1.8cm,yshift=-0.5cm]current page.north east);
        \node at (logo) {
\includegraphics[height=.1\textheight]{northeastern.eps}};
    \end{tikzpicture}
}
\makeatother


\setbeamertemplate{footline}{
    \hbox{%
    \begin{beamercolorbox}[wd=\paperwidth,ht=3ex,dp=1.5ex,leftskip=2ex,rightskip=2ex]{page footer}%
        \usebeamerfont{title in head/foot}%
        \hfill
        \begin{tblr}{
            width=.8\linewidth,
            colspec={X[l]X[c]X[r]}
        }
            \insertsectionhead{}
            &
            \ifx\insertsubsection\empty
            \else
            \insertsubsection{} 
            \fi
            &
            \insertframenumber{} / \inserttotalframenumber
        \end{tblr}
        \hfill{}
    \end{beamercolorbox}}%
}

\defbeamertemplate{section page}{my progressbar}{
  \centering
  \begin{minipage}{22em}
    \raggedright
    \usebeamercolor[fg]{section title}
    \usebeamerfont{section title}
    \insertsection\\[-1ex]
    \usebeamertemplate*{progress bar in section page}
    \par
    \ifx\insertsubsectionhead\@empty\else%
      \usebeamercolor[fg]{subsection title}%
      \usebeamerfont{subsection title}%
      \insertsubsectionhead
    \fi
  \end{minipage}
  \par
  \vspace{\baselineskip}
}
\setbeamertemplate{section page}[my progressbar]


\DeclareRobustCommand{\projectname}{\textsc{WebServ}\xspace}

% Note. For the workshop paper template, both \title{} and \workshoptitle{} are required, with the former indicating the paper title shown in the title and the latter indicating the workshop title displayed in the footnote. 
% \title{\projectname: Full-stack Unified Scalable Environment for Web Agent Reinforcement Learning}
% \title{
% \projectname: A Browser-Server Environment for Efficient Training of Reinforcement Learning-based Web Agents at Scale
% }

\title{
\projectname: A Full-Stack and RL-Ready Web Environment for Training Web Agents at Scale 
}


\UseTblrLibrary{booktabs}
\graphicspath{ {./figures/} }

\author{
% Yuxuan Lu, Ziyi Wang, Jing Huang, Hui Liu, Jiri Gesi, Yan Han, Shihan Fu, Tianqi Zheng, Xianfeng Tang, Chen Luo, Yisi Sang, Jin Lai, Dakuo Wang
\textbf{Yuxuan Lu\textsuperscript{1,2}},
\textbf{Ziyi Wang\textsuperscript{1,2}},
\textbf{Jing Huang\textsuperscript{1}},
\textbf{Hui Liu\textsuperscript{1}},
\textbf{Jiri Gesi\textsuperscript{1}},
\textbf{Yan Han\textsuperscript{1}},
\textbf{Shihan Fu\textsuperscript{1,2}},
\textbf{Tianqi Zheng\textsuperscript{1}},
\textbf{Xianfeng Tang\textsuperscript{1}},
\textbf{Chen Luo\textsuperscript{1}},
\textbf{Yisi Sang\textsuperscript{1}},
\textbf{Jin Lai\textsuperscript{1}},
\textbf{Dakuo Wang\textsuperscript{1,2}}}


\institute{\textsuperscript{1}Amazon.com, Inc.,
\textsuperscript{2}Northeastern University}
\date{Jan 14 2026}
\begin{document}

\maketitle

\begin{frame}{Motivation}
  \begin{outline}
    \1 Large Language Model (LLM) Web Agents have been developed and widely used
    \1 WebGPT\cite{nakanoWebGPTBrowserassistedQuestionanswering2022}, Claude Computer Use, ChatGPT Atlas, Step \cite{sodhiStePStackedLLM2024}, LASER \cite{maLASERLLMAgent2024}, WebAgent \cite{gurRealWorldWebAgentPlanning2023}, \dots
  \end{outline}
\end{frame}

\begin{frame}{Motivation}
  \begin{outline}
    \1 These agents have become a primising approach for web automation tasks \cite{zhouWebArenaRealisticWeb2024} and simulating user behaviors \cite{luPromptingNotAll2025}. 
    \1 However, existing Web Agents are mainly built with \textbf{prompting-based} and \textbf{behavior cloning-based} methods.
    \1 RLHF and RLVR have been proven to work in other fields and tasks \cite{ouyangTrainingLanguageModels2022,deepseek-aiDeepSeekR1IncentivizingReasoning2025}
    \1 RL Web Agents is limited by the absense of a full-stack, efficient and effective environment.
  \end{outline}
\end{frame}

\begin{frame}{Web Browsers are NOT RL-Ready}
  \begin{outline}
    \1 Web Browsers are NOT RL-Ready:
      \2 Requires per-site configuration or rely on site features
      \2 Noisy context and action
      \2 Lack of visual hint
      \2 Non-robust action execution
  \end{outline}
\end{frame}

\begin{frame}{Web Servers are NOT RL-Ready}
  \begin{outline}
    \1 Web Servers are usally packed in docker images
    \1 Not possible to run hundsreds \textit{isolated} web server containers to support parallel RL rollouts
  \end{outline}
\end{frame}

\begin{frame}{Existing Work}
  \begin{outline}
    \1 Researchers have argued that existing web environments are not RL-ready \cite{chenScalingAgentLearning2025}
    \1 Existing works are:
      \2 \dots off-policy \cite{qiWebRLTrainingLLM2025}
      \2 \dots not isolated \cite{weiWebAgentR1TrainingWeb2025}
      \2 \dots uses synthetic environments \cite{chenScalingAgentLearning2025}
    \1 The community needs a scalable and efficient training environment 
  \end{outline}
\end{frame}

\section{\projectname Architecture Design}

\begin{frame}{\projectname}
  \begin{figure}
    \includegraphics[width=.6\linewidth]{figures/teaser.png}
    \caption{System Architecture of \projectname compared to existing environments}
  \end{figure}
\end{frame}

\begin{frame}{\projectname: Web Browser Interface Design}
  \begin{outline}
    \1 To tackle challenges in web browser interface, we design:
      \2 Fully automatic parser 
      \2 Uses a simplified HTML as context
        \3 LLM already recognizes HTML structure
      \2 Network-Idle based waiting 
      \2 VLM Support
      \2 
  \end{outline}
\end{frame}

\begin{frame}{\projectname: Web Browser Interface Design}
  \begin{outline}
    \1 LLM Agent often click on non-interactive elements, but human can easily tell text and link apart
  \end{outline}
  \pause
  \begin{figure}
    \centering
    
\includegraphics[width=.4\linewidth]{image1.png}
    \caption{Visual Hint Example}
  \end{figure}
  \pause
  \begin{outline}
    \1 We made these visual hints available for agent as well
  \end{outline}
\end{frame}

\begin{frame}{\projectname: Web Server Design}
  \begin{outline}
    \1 Efficient Incus-based server container management
      \2 Run a server in seconds instead of minutes
    \1 Reduce resource need by 240x
    \1 Make it possible to run 200 parallel instances
  \end{outline}
\end{frame}

\section{Evaluation and Experiments}

\begin{frame}{Evaluation}
  \begin{outline}
    \1 Web browser I/O interface: Task successful rate on WebArena \cite{zhouWebArenaRealisticWeb2024}
    \1 Server manager effiency: Launch speed and resource need
  \end{outline}
\end{frame}

\begin{frame}{Evaluation Result}
  \begin{table}[t]
    \centering
    \begin{booktabs}{
      colspec = {X c c c},
      width=\linewidth,
      row{1} = {font=\bfseries},
      cell{2,7}{1} = {c=4}{bg=black!30!white,font=\bfseries},
    }
    \toprule
    Model and Method & Shopping & CMS & GitLab \\
    \midrule
    Vanilla WebArena \\
    % Qwen2.5-3B \cite{weiWebAgentR1TrainingWeb2025}  & 4.4 & 5.7 & 13.3 \\
    % Llama3.1-8B \cite{weiWebAgentR1TrainingWeb2025}  & 8.9 & 5.7 & 10.0 \\
    % Qwen2.5-32B \cite{weiWebAgentR1TrainingWeb2025}  & 17.8 & 20 & 20.0 \\
    GPT-4o  & 11.1 & 20 & 10.0 \\
    % GPT-4o-Turbo \cite{weiWebAgentR1TrainingWeb2025}  & 13.3 & 14.3 & 16.7 \\
    OpenAI-o3 & 33.3 & 45.7 & 46.7 \\
    % Step \cite{sodhiStePStackedLLM2024} & Multi-Prompt & 37.0 & 24.0 & 30.0\\
    % Llama3.1-8B + Behavior Cloning \cite{weiWebAgentR1TrainingWeb2025} & Fine-Tuning & 17.8 & 20 & 6.7 \\
    Llama-3.1-8B & 8.9 & 5.7 & 10.0 \\
    WebAgent-R1 & \textbf{44.4} & \textbf{57.1} & \textbf{56.7} \\
    % \midrule
    \projectname \\
    % Qwen-2.5-32B   & 17.8 & 20.0 & 20.0 \\
    GPT-4o    & 20.0 & 28.6 & 43.3 \\
    % GPT-4o Mini    & 20.0 & 25.7 & 30.0 \\
    GPT-5   & 35.6 & 57.1 & \textbf{53.3} \\
    Claude 3.5 Sonnet    & 26.7 & 31.4 & 36.7 \\
    Claude 3.7 Sonnet    & 31.1 & 37.1 & 50.0 \\
    Claude 4 Sonnet    & \textbf{42.2} & 48.6 & 50.0 \\
    Claude 4.5 Sonnet    & 40.0 & \textbf{62.9} & 50.0 \\
    % DeepSeek R1 & Reasoning Model & 13.3 & 37.1 & 23.3 \\
    \bottomrule
    \end{booktabs}
    % \vspace{2em}
    \caption{Comparison of model accuracy (in \%) across Shopping, CMS and GitLab tasks in WebArena-Lite.}
    \end{table}
\end{frame}

\begin{frame}{Evaluation Result}
  \begin{outline}
    \1 With \projectname, model can achieve better performancde than with vanilla WebArena
    \1 Achieves tate-of-the-art single-prompt agent performance
    \1 Prompting-based baseline models can achieve and even beat RL based methods
  \end{outline}
\end{frame}

\begin{frame}{Ablation Study: Visual Cue}
  \begin{outline}
    \1 How important are visual cues?
    \1 Ablation study: remove visual cues and "clickable" element hint in the observation space.
  \end{outline}
\end{frame}

\begin{frame}{Ablation: Visual Cue}
  \begin{table}[t]
    \centering
    \footnotesize
    \begin{booktabs}{
      colspec={Xccc},
      column{1}={l,font=\footnotesize},
      row{1}={font=\bfseries\footnotesize,c},
      width=\linewidth
    }
    \toprule
    Model & $\Delta$ Shopping (\%) & $\Delta$ CMS (\%) & $\Delta$ GitLab (\%) \\
    \midrule
    GPT-4o & \highlightRelevantValue{-33.3} & \highlightRelevantValue{-70.0} & \highlightRelevantValue{-46.2} \\
    GPT-5 & \highlightRelevantValue{-18.8} & \highlightRelevantValue{-35.0} & \highlightRelevantValue{-25.0} \\
    Claude 3.5 Sonnet & \highlightRelevantValue{-58.3} & \highlightRelevantValue{-45.5} & \highlightRelevantValue{-81.8} \\
    Claude 3.7 Sonnet & \highlightRelevantValue{-50.0} & \highlightRelevantValue{-53.8} & \highlightRelevantValue{-66.7} \\
    Claude 4 Sonnet & \highlightRelevantValue{-26.3} & \highlightRelevantValue{-29.4} & \highlightRelevantValue{-33.3} \\
    Claude 4.5 Sonnet & \highlightRelevantValue{-5.6} & \highlightRelevantValue{-45.5} & \highlightRelevantValue{6.7} \\
    \bottomrule
    \end{booktabs}
    \caption{Performance difference (\%) after removing visual cues.}
  \end{table}
  \begin{outline}
    \1 Except for one setting, all models exhibit performance drops
    \1 Weaker models suffer larger performance drops
    \1 Proves that Visual Cues is important especially for weaker models.
  \end{outline}
\end{frame}

\begin{frame}{Ablation: Vision Language Model}
  \begin{outline}
    \1 How important is other vision signals?
  \end{outline}


\begin{table}[t]
  \centering
  \begin{booktabs}{
  colspec={lccc},
  cell{3,5}{1}={r}
  }
  \toprule
  Model & Shopping & CMS & GitLab \\
  \midrule
  Claude 4 Sonnet & 42.2 & \textbf{48.6} & \textbf{50.0} \\
   + VLM & \textbf{44.4} & \textbf{48.6} & 43.3 \\
  Claude 4.5 Sonnet & \textbf{40.0} & 62.9 & 50.0 \\
   + VLM & 37.8 & \textbf{65.7} & \textbf{56.7} \\
  \bottomrule
  \end{booktabs}
  \caption{Ablation results comparing text-only and VLM settings for Claude~4 Sonnet and Claude~4.5 Sonnet across WebArena-Lite tasks.}
  \end{table}
  \pause
  Visual signals are not universally beneficial for general web agent tasks
\end{frame}

\begin{frame}{Evaluation: Launch Speed and Resource Need}
  \begin{table}[t]
    \centering
    \begin{booktabs}{
      colspec = {l c c},
      row{1} = {font=\bfseries}
    }
    \toprule
    Metric & \projectname (Incus) & Naïve Docker \\
    \midrule
    Launch speed & 1.781\,s & 8.963\,s \\
    Storage & 28.01\,MiB & 6.78\,GiB \\
    Memory & 1.74\,GiB & 1.63\,GiB \\
    \bottomrule
    \end{booktabs}
    \caption{Comparison of system efficiency between \projectname and Docker.}
    \label{tab:system_speed}
    \end{table}
    With Incus, \projectname can reduce launch latench and improved storage footprint with comparable memory usage, making large-scale RL rollouts possible. 
\end{frame}

\begin{frame}{Conclusion}
  \begin{outline}
    \0 WebServ enables:
      \1 State-of-the-art single-prompt agent performance on WebArena-Lite
      \1 No per-site configuration, able to operate on real complex website (Amazon, Google Flight, etc.)
      \1 Efficient RL training with VeRL integration
  \end{outline}
\end{frame}

\end{document}